\section*{Capítulo \ref{ch:b1:titration-methods} --- Métodos e curvas de titulação}

\begin{solution}{exo:b1:titration-methods:1}
\ce{H3O+ + OH- -> 2H2O}, $K = 1/K_e = 10^{14.00}$. $C = 12.4 \times
0.100/20.0 = \qty{0.0620}{mol/L}$.
\end{solution}

\begin{solution}{exo:b1:titration-methods:2}
Na \omterm{def:b1:titration-methods:titration}{equivalência}, a solução é de cloreto de amônio a
\qty{0.050}{mol/L}: $\mathrm{pH} = \frac12(9.25 - \log 0.050) = 5.28$.
O vermelho de metila (3.8--5.8) o contém; o alaranjado de metila, cuja faixa (2.5--4.5) só é atingida depois da \omterm{def:b1:titration-methods:titration}{equivalência}, quando o pH cai, mudaria tarde demais.
\end{solution}

\begin{solution}{exo:b1:titration-methods:3}
\ce{MnO4- + 5Fe^2+ + 8H+ -> Mn^2+ + 5Fe^3+ + 4H2O};
$C = 5 \times 0.0200 \times 12.6/10.0 = \qty{0.126}{mol/L}$.
\end{solution}

\begin{solution}{exo:b1:titration-methods:4}
$C = 0.0100 \times 14.2/50.0 = \qty{2.84e-3}{mol/L}$ de cálcio mais
magnésio. Os dois reagem antes do \omterm{def:b1:titration-methods:titration}{ponto final}, que só dá a sua soma:
uma \omterm{def:b1:titration-methods:kinds}{titulação simultânea}.
\end{solution}

\begin{solution}{exo:b1:titration-methods:5}
0~mL: \omterm{def:b1:predominant-reaction:strong-acid}{ácido fraco}, $\mathrm{pH} = \frac12(4.76 + 1.00) = 2.88$. 5.0~mL:
semiequivalência, 4.76. 10.0~mL: 8.73
(\cref{prop:b1:titration-methods:ph-at-equivalence}). 15.0~mL: excesso
de hidróxido de \qty{0.50}{mmol} em \qty{25.0}{mL}, $[\ce{OH-}] = 0.020$,
$\mathrm{pH} = 12.30$. Todos concordam com a curva calculada a 0.01.
\end{solution}

\begin{solution}{exo:b1:titration-methods:6}
Os $\mathrm{p}K_a$ 2.15 e 7.21 diferem de mais de 4, assim como 7.21 e
12.34: duas \omterm{def:b1:titration-methods:titration}{equivalências}, em 10.0 e \qty{20.0}{mL}. Primeira: solução
de \ce{H2PO4-}, $\mathrm{pH} \approx \frac12(2.15 + 7.21) = 4.68$;
segunda: \ce{HPO4^2-}, $\frac12(7.21 + 12.34) = 9.78$ (a curva exata dá
4.71 e 9.66, pois as fórmulas desprezam a diluição e a água). A terceira
acidez ($\mathrm{p}K_a$ 12.34) está próxima demais da da água: o
hidróxido adicionado só reage em parte e o pH sobe sem salto.
\end{solution}

\begin{solution}{exo:b1:titration-methods:7}
Ácido clorídrico: $6.2 \times 0.100/10.0 = \qty{0.062}{mol/L}$; ácido
etanoico: $(14.6 - 6.2) \times 0.100/10.0 = \qty{0.084}{mol/L}$. A meio
caminho, metade do ácido etanoico está titulada: $\mathrm{pH} = \mathrm{p}K_a = 4.76$.
\end{solution}

\begin{solution}{exo:b1:titration-methods:8}
$V_{\mathrm{eq}} = \qty{10.0}{mL}$. 2.0~mL: $E = 0.77 + 0.059\log(2/8) =
0.73$~V; 9.0~mL: $0.77 + 0.059\log 9 = 0.83$~V; 11.0~mL: excesso de
permanganato de \qty{0.020}{mmol}, \ce{Mn^2+} \qty{0.200}{mmol}, $E = 1.51 +
\frac{0.059}{5}\log 0.10 = 1.50$~V; 20.0~mL: 1.51~V.
\end{solution}

\begin{solution}{exo:b1:titration-methods:9}
0~mL: $\kappa = (349.8 + 76.2) \times 0.0100 = \qty{4.26}{mS/cm}$.
10.0~mL: \ce{Na+} e \ce{Cl-} a $1.00/110 = \qty{9.09e-3}{mol/L}$,
$\kappa = (50.3 + 76.2) \times \num{9.09e-3} = 1.15$. 20.0~mL: \ce{Na+}
$0.0167$, \ce{Cl-} e \ce{OH-} $0.00833$~mol/L, $\kappa = 0.84 + 0.635 + 1.65
= 3.13$~mS/cm. Inclinações (diluição desprezada, \qty{1}{mL} acrescenta
\qty{1.0e-3}{mol/L}): $(50.3 - 349.8) \times 10^{-3} = -0.30$ e
$(50.3 + 198.3) \times 10^{-3} = +0.25$~mS/cm por mL.
\end{solution}

\begin{solution}{exo:b1:titration-methods:10}
Em $V_{\mathrm{eq}} - v$, o excesso de ácido é $C_Bv$; em $V_{\mathrm{eq}} + v$,
o excesso de base é $C_Bv$; no mesmo volume, $[\ce{H3O+}]$ antes é igual
a $[\ce{OH-}]$ depois, logo $\mathrm{pH}(V_{\mathrm{eq}} - v) +
\mathrm{pH}(V_{\mathrm{eq}} + v) = 14$: simetria em relação a
(V$_{\mathrm{eq}}$, 7). Salto: \qty{1.0e-5}{mol} de excesso em
\qty{20}{mL}, pH de 3.30 a 10.70, 7.4 unidades. Dividindo por 100:
$\num{1.0e-7}$~mol em \qty{20}{mL}, pH de 5.3 a 8.7, 3.4 unidades: um
\omterm{def:b1:titration-methods:titration}{ponto final} muito menos nítido.
\end{solution}

\begin{solution}{exo:b1:titration-methods:11}
Antes: \ce{Fe^3+} formado $= 5C_{\mathrm{Mn}}V$, \ce{Fe^2+} restante $=
5C_{\mathrm{Mn}}(V_{\mathrm{eq}} - V)$, logo $E = 0.77 + 0.059\log\frac{V}{V_{\mathrm{eq}}
- V}$: em $V_{\mathrm{eq}}/2$, $E = 0.77$. Depois: excesso de \ce{MnO4-} $\propto
V - V_{\mathrm{eq}}$, \ce{Mn^2+} $\propto V_{\mathrm{eq}}$: $E = 1.51 +
\frac{0.059}{5}\log\frac{V - V_{\mathrm{eq}}}{V_{\mathrm{eq}}}$ (pH 0); em
$2V_{\mathrm{eq}}$, 1.51~V. Na \omterm{def:b1:titration-methods:titration}{equivalência}, $(0.77 + 5 \times 1.51)/6 =
1.39$~V. Em pH 1, o par do permanganato fica mais baixo de $0.059 \times 8/5 =
0.094$~V: $E_{\mathrm{eq}} = (0.77 + 5 \times 1.416)/6 = 1.31$~V.
\end{solution}

\begin{solution}{exo:b1:titration-methods:12}
\ce{NH4+ + OH- -> NH3 + H2O} tem $K = 10^{14.00 - 9.25} = 10^{4.75}$: a
reação não é completa perto da \omterm{def:b1:titration-methods:titration}{equivalência} e o salto é pequeno, em
torno de pH 11 (pH de \omterm{def:b1:titration-methods:titration}{equivalência} $14 - \frac12(4.75 + 1.30) = 11.0$).
\omterm{def:b1:titration-methods:kinds}{Titulação de retorno}: adicione $n_1$ de hidróxido de sódio (excesso),
ferva para expulsar a amônia formada e titule o hidróxido restante com
ácido clorídrico ($n_2$, salto nítido forte--forte): $n(\ce{NH4+}) = n_1 - n_2$.
\end{solution}

\begin{solution}{pb:b1:titration-methods:1}
\textbf{1.} \ce{C6H6O6 + 2H+ + 2e- -> C6H8O6}.
\textbf{2.} \ce{C6H8O6 + I2 -> C6H6O6 + 2I- + 2H+}.
\textbf{3.} 0 e $-$I.
\textbf{4.} \ce{I2 + 2S2O3^2- -> 2I- + S4O6^2-}, $\log K = 2(0.53 - 0.02)/0.059
= 17$.
\textbf{5.} A cor azul exige iodo; ela desaparece com o último iodo, na
\omterm{def:b1:titration-methods:titration}{equivalência}.
\textbf{6.} As perdas para o ar diminuiriam o resultado; adicionado de
uma vez e em excesso, o iodo oxida o ácido ascórbico de modo rápido e
completo.
\textbf{7.} Uma reação rápida e \omterm{def:b1:extent-q-and-k:quantitative}{quantitativa} durante toda a \omterm{def:b1:titration-methods:titration}{titulação},
sem oxidação pelo ar enquanto o \omterm{def:b1:titration-methods:titration}{titulante} é adicionado lentamente, e um
\omterm{def:b1:titration-methods:titration}{ponto final} no primeiro excesso de iodo.
\textbf{8.} Excesso de iodo $n_R$; reação com o ácido ascórbico;
\omterm{def:b1:titration-methods:titration}{titulação} do iodo restante pelo tiossulfato.
\textbf{9.} Caso contrário, restaria ácido ascórbico e o excesso de iodo,
nulo, não diria quanto. Aqui reagiram
\qty{2.775e-4}{mol} de \qty{5.000e-4}{mol}: há de fato excesso.
\textbf{10.} As soluções de tiossulfato mudam lentamente com o tempo; sua
concentração entra diretamente no resultado.
\textbf{11.} 10.
\textbf{12.} Uma pipeta volumétrica de \qty{20.00}{mL}.
\textbf{13.} $20.00 \times 10^{-3} \times 0.0250 = \qty{5.000e-4}{mol}$.
\textbf{14.} $8.90 \times 10^{-3} \times 0.0500 = \qty{4.450e-4}{mol}$.
\textbf{15.} A metade: \qty{2.225e-4}{mol}.
\textbf{16.} $5.000 - 2.225 = \qty{2.775e-4}{mol}$.
\textbf{17.} Um para um: \qty{2.775e-4}{mol}.
\textbf{18.} $\times 10$: \qty{2.775e-3}{mol}, $\times 176.0$:
\qty{0.488}{g}.
\textbf{19.} Cerca de \qty{2}{\percent} abaixo do rótulo.
\textbf{20.} $n = \bigl(C_IV_I - \frac12C_TV_T\bigr)\,V_{\text{balão}}/V_{\text{alíquota}}$.
\textbf{21.} Iodo introduzido: pipeta $0.03/20.00 = \qty{0.15}{\percent}$
e concentração \qty{0.2}{\percent}, cerca de \qty{0.25}{\percent}, ou
seja, \qty{1.3e-6}{mol}. Excesso: bureta $0.05/8.90 = \qty{0.56}{\percent}$
e concentração \qty{0.2}{\percent}, cerca de \qty{0.60}{\percent}, ou
seja, \qty{1.3e-6}{mol}.
\textbf{22.} As incertezas absolutas se combinam, enquanto a diferença é
menor que cada uma das quantidades: $\sqrt{1.3^2 + 1.3^2} \times 10^{-6} =
\qty{1.8e-6}{mol}$, \qty{0.66}{\percent} de \qty{2.775e-4}{mol}, mais que
cada um dos termos.
\textbf{23.} Com o balão (\qty{0.1}{\percent}) e a pipeta da alíquota
(\qty{0.2}{\percent}): $\sqrt{0.66^2 + 0.1^2 + 0.2^2} = \qty{0.70}{\percent}$.
\textbf{24.} $0.488 \pm 0.003$~g: \textbf{\qty{0.49}{g} de ácido
ascórbico por comprimido}.
\end{solution}
